\documentclass[11pt,letterpaper]{article}
\usepackage[T1]{fontenc}
\usepackage[utf8]{inputenc}
\usepackage{amsmath,amsthm}
\usepackage{newtxtext,newtxmath}
\usepackage[margin=0.86in,headheight=15pt]{geometry}
\usepackage{microtype}
\usepackage{mathtools}
\usepackage{booktabs,longtable,array,enumitem}
\usepackage{xcolor,listings,xurl}
\usepackage{fancyhdr}
\usepackage[hidelinks,pdfusetitle]{hyperref}
\usepackage{bookmark}
\definecolor{ink}{HTML}{17344B}
\definecolor{muted}{HTML}{50606B}
\definecolor{codebg}{HTML}{F5F6F7}
\definecolor{coderule}{HTML}{D9DFE3}
\newcommand{\code}[1]{\texttt{\detokenize{#1}}}
\newcommand{\src}[1]{\textnormal{(source lines #1)}}
\newcommand{\Q}{\mathbb{Q}}
\newcommand{\C}{\mathbb{C}}
\newcommand{\Z}{\mathbb{Z}}
\DeclareMathOperator{\Arg}{Arg}
\DeclareMathOperator{\lcm}{lcm}
\DeclareMathOperator{\disc}{disc}
\DeclareMathOperator{\gcdop}{gcd}
\DeclareMathOperator{\Repart}{Re}
\DeclareMathOperator{\Impart}{Im}
\newtheorem{proposition}{Proposition}[section]
\newtheorem{lemma}[proposition]{Lemma}
\theoremstyle{definition}
\newtheorem{example}[proposition]{Example}
\newtheorem{definition}[proposition]{Definition}
\lstdefinelanguage{AuditWL}{
 morekeywords={Block,Module,If,For,While,Which,Return,ClearAll,Function,True,False,
   Null,Power,Times,Plus,List,Rule,RuleDelayed,Set,SetDelayed,MinimalPolynomial,
   PolynomialGCD,PolynomialQuotient,RootReduce,FactorList,FactorInteger,Sort,
   SortBy,Select,DeleteDuplicates,ReplaceAll,ReplaceRepeated,Part,TrueQ,Check,
   Quiet,TimeConstrained,Failure,Association,BeginPackage,Begin,End,EndPackage},
 sensitive=true,morecomment=[s]{(*}{*)},morestring=[b]"
}
\lstset{language=AuditWL,basicstyle=\ttfamily\footnotesize,
 keywordstyle=\color{ink},commentstyle=\color{muted},stringstyle=\color{muted},
 backgroundcolor=\color{codebg},frame=single,rulecolor=\color{coderule},
 columns=fullflexible,keepspaces=true,showstringspaces=false,
 breaklines=true,breakatwhitespace=false,tabsize=2,
 aboveskip=9pt,belowskip=9pt,numberstyle=\tiny\color{muted},numbersep=8pt,
 xleftmargin=4pt,xrightmargin=4pt}
\setlength{\parskip}{4.2pt}
\setlength{\parindent}{0pt}
\setlist{itemsep=2pt,topsep=5pt,leftmargin=1.5em}
\setlength{\emergencystretch}{2em}
\pagestyle{fancy}
\fancyhf{}
\fancyhead[L]{\small\color{muted}Radical denesting: technical audit}
\fancyhead[R]{\small\color{muted}5 September 2026}
\fancyfoot[C]{\small\thepage}
\renewcommand{\headrulewidth}{0.3pt}
\hypersetup{pdftitle={Technical Audit of a Wolfram Language Radical-Denesting Program},
 pdfauthor={OpenAI},pdfsubject={Source-level analysis, algebraic correctness, branch defects, and regression tests}}
\title{\vspace{-1.2em}\textbf{A Technical Audit of a Wolfram Language\\Radical-Denesting Program}\\[0.5em]
 \large Algebraic design, branch correctness, search heuristics,\\and engineering risks}
\author{Source-grounded analysis of the supplied \texttt{Pasted text.txt}}
\date{5 September 2026}
\begin{document}
\maketitle
\thispagestyle{empty}
\begin{abstract}
The supplied program attempts to simplify nested algebraic radicals by factoring
expressions, collecting radical subproblems, and searching for multiplicative
changes of variables that make a polynomial greatest common divisor have lower
degree. Its central algebraic idea is useful: the minimal polynomial of a suitable
multiple of a radical can isolate a simpler polynomial relation over an algebraic
coefficient field. This article reconstructs that mechanism, proves the associated
candidate-generation and denominator-rationalization identities, and follows the
implementation from its public wrappers to its auxiliary tree-processing routines.

The most consequential defect is a loss of the original complex branch: the code
replaces a target $\alpha$ by the principal root of $\alpha^k$ when selecting an
answer. Explicit examples show a sign reversal for a negative real target, a
complex $3/2$ power, and a product of two principal square roots. Separate defects
include branch-unsafe exponent manipulation in the custom factorizer, collisions
between temporary symbols and user expressions, an incorrect sorting comparator,
an unguarded access to an empty candidate list, and a supposed multiplier cap
that can allow billions of candidates. Other limitations are heuristic rather
than correctness defects: incomplete factorization, degree-based pruning,
unproven search coverage, and the absence of a verified denesting objective.

The audit is based on the complete supplied source, official Wolfram documentation,
and independently executed symbolic and numerical checks. The uploaded Wolfram
program and the proposed Wolfram regression tests were \emph{not} executed in a
Wolfram kernel. This distinction is maintained throughout.
\end{abstract}
\vfill
\begin{center}\small
\textbf{Source fingerprint (SHA-256)}\\[3pt]
\texttt{693a2825f1e2b8912e55371b90dc8ed115bc}\\
\texttt{33ee707dd3a4521eb22bcd2ccc93}\\[8pt]
The archive contains the article, its LaTeX source, the original code,\\
independent verification scripts and results, and proposed regression tests.
\end{center}
\newpage
\tableofcontents
\newpage

\section{Scope, evidence, and the intended contract}
\label{sec:scope}

\subsection{What was actually supplied}
The attachment is a Wolfram Language program, not a prose specification. It defines
\code{Strad}, \code{DenestRadicals3}, \code{denest11}, several mathematical helpers,
and a custom factorization subsystem. The original byte stream is preserved in
\code{original.wl}; its 585 physical lines are reproduced in Appendix~\ref{app:source}.
All source-line references in this article refer to that file. The attachment
provides no author attribution, version requirement, test suite, benchmark corpus,
formal input contract, or completeness theorem. None is inferred here.

The public-looking entry points are
\begin{lstlisting}
Strad[expr, alllevels, func]
DenestRadicals3[expr, alllevels : False, func : denest11]
denest11[problem, forcedmultipliers : 0]
\end{lstlisting}
The first signature above describes forwarding, not a separately declared
three-argument definition: the actual \code{Strad[eaf__]} accepts a nonempty
sequence and passes it onward \src{1--5}. The \code{forcedmultipliers} argument
belongs to \code{denest11}; it is not the second argument of the wrapper.

The wrapper recognizes powers whose exponents are exact \code{Rational} objects
and whose bases satisfy \code{AlgebraicQ}. The implementation of that predicate is
\code{Element[num, Algebraics]} \src{453--454}. The code therefore appears intended
primarily for explicit, exact algebraic numbers written using radicals. However,
its signatures do not enforce that domain, and \code{Factorc[expr]} is applied to
the entire input \emph{before} any algebraicity test. This matters for the symbolic
counterexamples below.

\subsection{What counts as a correct denester}
For an exact scalar algebraic input $\alpha$, a successful simplifier should at
least return an expression $\beta$ satisfying
\begin{equation}\label{eq:contract}
 \beta=\alpha
\end{equation}
under the language's branch conventions. A \emph{denesting} success should additionally
improve an explicitly chosen representation criterion. Equality alone is not a
proof of denesting; a lower minimal-polynomial degree for a transformed auxiliary
number is not a proof of denesting either.

For example, the same algebraic number can be written as a nested radical, a sum
of unnested radicals, or a \code{Root} object. A conversion to a \code{Root} object
can eliminate all visible \code{Power} nodes without providing an unnested radical
formula. Any practical contract must state which of these output forms are allowed.

\subsection{Evidence levels and execution boundary}
\begin{longtable}{@{}p{0.16\textwidth}p{0.80\textwidth}@{}}
\toprule
Evidence & Meaning in this article\\\midrule
\endhead
Source fact & A statement directly established by the supplied definitions, indices,
conditions, or assignments.\\
Algebraic proof & A mathematical consequence proved here independently of a particular
Wolfram kernel's evaluation choices.\\
Executed check & A computation actually run with SymPy 1.14.0 or mpmath 1.3.0.
These checks do not execute or emulate the uploaded program.\\
Predicted trace & A Wolfram-level behavior deduced from the source and documented
semantics; the corresponding regression test is supplied but was not run here.\\
Conditional risk & A failure mode whose occurrence depends on an input, a callback,
resource exhaustion, or solver behavior not established by a complete runtime trace.\\
\bottomrule
\end{longtable}
No local Wolfram executable was available, and no connected Wolfram computation
service was used. Consequently, this report does not claim a Mathematica version,
end-to-end timing, an observed Wolfram message transcript, or a passing Wolfram test
count. Some concrete examples can be simplified automatically before reaching a
faulty path in particular releases. Their mathematical role is to expose the
invalid transformation; the regression suite determines runtime manifestation.

Official documentation is used only to establish language semantics and built-in
function behavior. The algorithmic interpretation, counterexamples, proofs, and
recommendations are the analysis developed from the supplied source.

\section{Architecture: from an expression to a multiplier search}
\label{sec:architecture}

\begin{longtable}{@{}>{\raggedright\arraybackslash}p{0.30\textwidth}p{0.14\textwidth}>{\raggedright\arraybackslash}p{0.51\textwidth}@{}}
\toprule
Routine & Lines & Responsibility\\\midrule
\endhead
\code{Strad} & 1--3 & Forward arguments while suppressing calls to \code{Print}.\\
\code{DenestRadicals3} & 4--76 & Preprocess the expression, mark radicals, combine
subproblems, dispatch to a denester, and substitute results.\\
\code{denest11} & 78--441 & Choose a root order, generate multipliers, manage search
history and a mutable worklist, and retain a proposed solution.\\
\code{trymultiplier} & 107--155 & Compute a minimal polynomial and possibly a
lower-degree GCD; solve and select a candidate. Local to \code{denest11}.\\
\code{newmultipliers} & 159--211 & Generate additional multipliers using discriminant
factors or a product of previously useful multipliers. Also local.\\
\code{RationalizeDenominator} & 443--451 & Replace the reciprocal of an algebraic
denominator by a polynomial in that denominator.\\
\code{AlgebraicQ} & 453--454 & Ask whether an expression belongs to the algebraic numbers.\\
\code{nestdepths} & 456--466 & Recursively annotate an expression tree with a custom
nesting height.\\
\code{nestdepthsort} & 468--480 & Flatten and group those annotations by height.\\
\code{maxnestdepth} & 482--498 & Extract the relevant maximum-height frontier.\\
\code{ComplexitySort} & 500--504 & Sort first by \code{LeafCount}, then by magnitude.\\
\code{Factorc}, \code{fullfactor}, \code{combine} & 506--584 & Encode products and sums
as nested lists, extract common factors, and reconstruct an expression.\\
\bottomrule
\end{longtable}

The principal flow is
\begin{center}
\code{Strad} $\longrightarrow$ \code{DenestRadicals3}
$\longrightarrow$ \code{denest11}\\[4pt]
$\hspace{1.5em}\searrow$ factorization and tree helpers
$\hspace{1em}\searrow$ multiplier trials and new search frames.
\end{center}
The auxiliary factorizer is not a harmless formatting pass. It changes the actual
expression that every later stage sees. A correctness check performed only inside
\code{denest11} cannot repair a value already changed by this preprocessing.

\section{The wrapper in detail}
\label{sec:wrapper}

\subsection{Marking and preprocessing radical subexpressions}
The value of \code{alllevels} chooses \code{ReplaceAll} unless it is literally
\code{True}; only that exact value chooses \code{ReplaceRepeated} \src{7--10}.
Thus a typo or non-Boolean value silently selects the default mode rather than
producing an option error.

For each matching rational power, the wrapper splits its base into a numerical
content and a remaining factor using \code{FactorTermsList}. It further processes
the remainder through \code{Factorc}. The numerical content is decomposed using
the GCD of its real and imaginary parts and then \code{FactorInteger}
\src{16--25}. The local predicate \code{SumQ} is not merely a test of whether the
outer head is \code{Plus}: it searches for a sum at any level and also recognizes
an explicit complex number with nonzero real part \src{13--15}.

Factors identified by this predicate are multiplied together and passed through
\code{ComplexExpand}. Factors with numerically negative real part are negated,
while a separate sign accumulator records the compensating sign. Finally, the
factors are rationalized individually, multiplied back together, and simplified
\src{26--40}.

There is an important distinction here. The sign accumulator preserves the product
algebraically: replacing a factor $a$ by $-a$ while multiplying the accumulated
sign by $-1$ does not change the product. An inaccurate sign classification can
make the heuristic less useful, but it is not by itself the branch-loss defect
proved later. Similarly, \code{ComplexExpand} assumes unspecified variables are
real, but on the intended explicit algebraic-number domain there should be no such
variables to assume real \cite{complexexpand}. It is therefore inappropriate to
attribute every complex-branch problem to that built-in.

The processed radical becomes an internal object \code{f[base, exponent]}. The
symbol \code{f} is localized by \code{Block}, not made fresh by \code{Module};
Section~\ref{sec:collisions} explains why this distinction matters.

\subsection{Combining subproblems}
Two marked factors are combined into one marked \emph{product of their powers}
when both bases contain a \code{Power} node \src{43--47}. This is not the naive
identity $(ab)^r=a^rb^r$: the code actually retains the product
\code{(rad1^pow1)*(rad2^pow2)} inside the marker. The grouping itself therefore
preserves the value.

Next, two-argument markers are changed into one-argument markers containing the
powered expression \src{48}. Another rule combines marked complex values when
both real and imaginary parts are numerically nonzero at precision 10
\src{49--54}. Again, the product is retained. The danger is downstream: the
combined product need not be a principal $k$th root, while the engine later treats
it as one.

\subsection{Single-pass versus all-level processing}
The wrapper collects distinct occurrences of \code{f[_]} throughout the expression
and logs them \src{55--57}. In default mode, it constructs a substitution rule for
each collected problem using the selected callback \code{func}, then simplifies
the result \src{58--61}.

In all-level mode, inner markers may occur inside outer markers. The code cycles
through the problem list until it finds a marker with no unresolved marker inside
it. It replaces that entry by a rule, propagates the rule through the remaining
problems, and eventually applies all rules to the main expression \src{62--73}.
This is a manually implemented bottom-up dependency resolution.

For a finite, well-formed marker tree and a callback that returns marker-free
expressions, the intended process can make progress by resolving an innermost
marker. It is not accurate to call this loop intrinsically nonterminating.
Nevertheless, it has no explicit check that a full cycle resolved something,
no callback-result contract, and no protection against a user expression that
already uses the same marker symbol.

The meaning of \code{alllevels -> True} is thus structural: process the marked
levels of the supplied expression. It does not establish that all mathematically
possible denestings are found, nor does it prove that every newly introduced
radical will itself be denested to a fixed point.

\section{The algebra behind the engine}
\label{sec:algebra}

\subsection{Choice of a root order}
After a syntactic nested-power test, \code{denest11} asks \code{maxnestdepth} for
rational powers at the maximum relevant nesting height. It sets
\begin{equation}\label{eq:rootorder}
 k=\lcm\{\text{denominators of those exponents}\},\qquad
 \texttt{outerpower}=\frac1k,\qquad A=\alpha^k,
\end{equation}
where $\alpha$ is the original \code{problem} \src{84--94}. The intended effect is
to remove the outer radical layer. This $k$ is selected from a syntactic frontier;
it is not a proven minimal field-theoretic root order, nor is it necessarily the
LCM of every rational exponent anywhere in the expression.

A second quantity, \code{extrapower}, is calculated from deeper levels and printed
\src{95--103}. It is never subsequently used to change the target, the multiplier
search, or candidate validation. The program therefore does not implement an
additional lifting step merely because it prints an ``extra root.''

\subsection{A trial multiplier}
For a nonzero algebraic multiplier $m$, a trial constructs
\begin{equation}\label{eq:beta}
 \beta=(mA)^{1/k},\qquad P_m(X)=\operatorname{minpoly}_{\Q}(\beta).
\end{equation}
Here and below rational powers have their principal complex meaning.
\code{MinimalPolynomial[s]} legitimately returns a pure-function representation;
calling that result on \code{arg1} is supported syntax, not a missing polynomial
variable bug \cite{minpoly}. The code records $d_m=\deg P_m$ \src{111--114}.

For selected trials, it computes
\begin{equation}\label{eq:gcd}
 G_m(X)=\gcd_{K[X]}\bigl(P_m(X),X^k-mA\bigr),
\end{equation}
using \code{Extension -> Automatic} \src{115--120}. The field $K$ contains the
algebraic numbers present in the polynomial coefficients. The representation of
those coefficients can expose a larger field than the smallest field generated
by their evaluated values; it is safer to describe $K$ this way than to assert
that it is always exactly $\Q(mA)$. The documented purpose of the option is to
include algebraic coefficients in the coefficient field \cite{polygcd}.

\begin{proposition}[Why the GCD is meaningful]
Assume $A,m\ne0$ are algebraic, $k\ge2$, and the operations are exact. Then
$1\le\deg G_m\le k$. Every root $t$ of $G_m$ satisfies $t^k=mA$.
\end{proposition}
\begin{proof}
The chosen value $\beta$ is a common root of the two polynomials in
\eqref{eq:gcd}. Its minimal polynomial over $K$ divides both, so their GCD is
nonconstant. The GCD divides $X^k-mA$, giving the degree bound and the asserted
root equation.
\end{proof}

When $\deg G_m<k$, the code regards the trial as having found a denesting. A
linear GCD is especially useful: it directly expresses one root in the coefficient
field. A quadratic or cubic GCD may provide a smaller equation, but its solutions
still need to be inspected to determine whether they improve the requested
radical representation.

\subsection{Why the root-of-unity enumeration is a good idea}
The code solves $G_m(t)=0$, divides the resulting roots by the principal
$m^{1/k}$, and multiplies by every $k$th root of unity \src{121--124}.

\begin{proposition}[Candidate coverage for a fixed successful GCD]
Let $t$ be any root of $G_m$ and $m\ne0$. Then the set
\begin{equation}\label{eq:orbit}
 \left\{\frac{t}{m^{1/k}}\exp\!\left(\frac{2\pi i j}{k}\right):
 j=0,1,\ldots,k-1\right\}
\end{equation}
is exactly the set of $k$th roots of $A$. In particular, it contains the original
$\alpha$ whenever $A=\alpha^k$.
\end{proposition}
\begin{proof}
Since $(m^{1/k})^k=m$, the $k$th power of $t/m^{1/k}$ is $A$. Multiplication by
the $k$ roots of unity gives all roots of $X^k-A$, which are distinct because
$A\ne0$ and the characteristic is zero.
\end{proof}

This argument does \emph{not} require the generally false identity
$(mA)^{1/k}=m^{1/k}A^{1/k}$. Enumerating the root-of-unity orbit is precisely what
can compensate for the choice of principal roots. The implementation already
creates the right type of candidate set. Its major mistake is selecting from
that set against $A^{1/k}$ instead of against the original $\alpha$.

\section{A complete worked example}
\label{sec:worked}
Consider
\begin{equation}
 \alpha=\sqrt{5+2\sqrt6}=\sqrt2+\sqrt3.
\end{equation}
Here $k=2$ and $A=5+2\sqrt6$. At $m=1$,
\begin{equation}
 P_1(X)=X^4-10X^2+1.
\end{equation}
Over $K=\Q(\sqrt6)$,
\begin{equation}
 G_1(X)=X^2-5-2\sqrt6.
\end{equation}
The GCD still has degree $k=2$, so this trial does not yield a smaller equation.
It does provide the first finite degree record for the search.

At $m=2$, however,
\begin{equation}
 \beta=\sqrt{10+4\sqrt6}=2+\sqrt6,
\end{equation}
and therefore
\begin{align}
 P_2(X)&=X^2-4X-2,\\
 G_2(X)&=X-2-\sqrt6.
\end{align}
The selected root gives
\begin{equation}
 \frac{2+\sqrt6}{\sqrt2}=\sqrt2+\sqrt3.
\end{equation}
All roots and signs agree with the original positive target in this example.

This illustrates why a multiplier helps: it moves the transformed root into a
smaller absolute-degree situation and, more importantly, makes the GCD linear over
the available coefficient field. The degree drop is a useful clue here, but the
linear GCD is the actual algebraic information used to construct the expression.

The discriminant of $P_1$ is
\begin{equation}
 \disc(P_1)=147456=2^{14}3^2.
\end{equation}
Thus the discriminant-based generator can discover prime multipliers 2 and 3 in
this example. This successful illustration is not a theorem that discriminant
primes suffice for arbitrary inputs.

The displayed minimal polynomials, GCDs, discriminant, and identity were checked
independently with SymPy. The order in which the uploaded program reaches $m=2$
was not measured in a Wolfram execution.

\section{Search state, multiplier generation, and stopping}
\label{sec:search}

\subsection{Initial multipliers}
When \code{forcedmultipliers} is a list, the code uses that list to seed the search.
Otherwise it expands a rationalized version of $A$, removes selected rational
coefficients from the putative terms, and forms complementary powers of radicals
\src{214--249}. For a term $b^{u/v}$ with $u\ne0$, the intended local calculation is
\begin{equation}
 \left(b^{u/v}\right)^{(v-u)/u}\approx b^{(v-u)/v},
\end{equation}
so that multiplication by the original term would remove its fractional exponent
at a formal level. Complex branch laws mean this mnemonic is not an unrestricted
identity. Because these expressions are merely trial multipliers, their equality
to some ``ideal'' multiplier is not required for soundness, provided the multiplier
is nonzero and the eventual result is certified against the original input.

The list starts with multiplier 1 in automatic mode. A list passed as
\code{forcedmultipliers} does not strictly prohibit later generated multipliers:
progress-triggered insertion is still present in the main loop. Its practical
meaning is closer to ``initial multipliers'' than to ``the only multipliers that
may be tried.''

\subsection{History and worklist representation}
The history starts with a sentinel
\begin{lstlisting}
{0, Infinity, Infinity}
\end{lstlisting}
and records triples \code{{multiplier, minpoly, degree}}. The history insertion
logic orders records by decreasing degree, placing the smallest degree at the
end \src{380--388}. The sentinel makes the first finite-degree trial count as
progress without a separate initialization branch.

The so-called \code{multiplierstack} is an ordered, mutable worklist rather than
an ordinary last-in-first-out stack. A realized frame has the shape
\begin{lstlisting}
{{doFullSet, generatingAttempt}, multiplierList}
\end{lstlisting}
and a deferred frame stores a function and its arguments until it is reached
\src{253--267, 346--348}. During processing, entries in a multiplier list change
from scalar multipliers into pairs \code{{multiplier, degree}} \src{292}.

If the current frame is not a full-set frame, a newly generated frame is inserted
immediately after the completed prefix and before any remaining suffix. In a
full-set frame, additional work is scheduled later, using degree and expression
size as ordering clues \src{350--375}. This is a reasonable search-management
idea, but the same nested-list structure carries several different schemas.
Its invariants are implicit, and no shape validation protects the many positional
accesses.

\subsection{New multipliers}
In the usual branch, the generator takes a partial factorization of the
discriminant of a recorded minimal polynomial, discards factors that do not pass
\code{PrimeQ}, and uses selected primes to form divisors of
\begin{equation}\label{eq:divisorgrid}
 \left(\prod_{j=1}^{r}p_j\right)^{k-1}.
\end{equation}
These divisors, together with additional discovered primes, are multiplied by the
previous multiplier \src{169--199}.

In the alternative branch, where observed degrees have a GCD smaller than either
one, the generator tries the product of the two recorded multipliers
\src{201--207}. The code supplies no field-theoretic proof that this product will
realize the numerical GCD of the degrees. It should be read as a search heuristic.

After the first automatic pass, the program also builds products of subsets of
some promising initial multipliers, then reduces certain exponent patterns
\src{396--426}. Several transformations in this block are not general complex
identities. Again, that is not automatically a returned-value bug: arbitrary
nonzero algebraic multipliers are legitimate search proposals. It does limit the
meaning of any claim that the generated list represents all equivalent cases.

\subsection{When the code declares success}
An attempt with a non-list last component is interpreted as carrying a solution.
The solution is simplified and sometimes returned immediately, depending on its
absolute degree relative to an earlier record \src{309--324}. This result protocol
uses the shape of the last element as a type tag. It assumes that a valid solution
is scalar and that no malformed expression has accidentally occupied that slot.

The program does not compare the new expression's radical nesting height with
that of the input before declaring success. Later ``additional denesting'' attempts
can replace an earlier solution without an explicit objective proving that the
replacement is a better radical expression. \code{Simplify} is used, but no
custom denesting complexity function is supplied.

\section{The auxiliary mathematics and data structures}
\label{sec:helpers}

\subsection{Denominator rationalization: a valid idea}
Write an input as $N/d$ with nonzero algebraic denominator $d$, and let $P$ be its
minimal polynomial. Polynomial division over a field containing $d$ gives
\begin{equation}
 P(X)=(X-d)Q(X),
\end{equation}
since $P(d)=0$. Evaluation at zero yields
\begin{equation}\label{eq:rationalization}
 P(0)=-dQ(0),\qquad \frac1d=-\frac{Q(0)}{P(0)}.
\end{equation}
Because $d\ne0$, its irreducible minimal polynomial cannot have zero constant term.
The code implements exactly this construction \src{443--451}.

The argument does not require a monic polynomial. This is relevant because
\code{MinimalPolynomial} uses a primitive integer polynomial with positive
leading coefficient, which may be nonmonic \cite{minpoly}. For example,
\begin{align}
 \frac1{1+\sqrt2}&=\sqrt2-1,\\
 \frac1{\sqrt2/3}&=\frac{3\sqrt2}{2},\\
 \frac1{1+i\sqrt2}&=\frac{1-i\sqrt2}{3}.
\end{align}
All three identities, including the nonmonic case, were independently checked.
The variable called \code{conjugate} actually holds a minimal-polynomial function,
not a single algebraic conjugate.

The mathematical formula is sound on its stated domain. The implementation still
has a temporary-symbol collision discussed below and lacks explicit handling of
zero, symbolic or nonalgebraic denominators, failed polynomial operations, and
resource limits.

\subsection{Nesting-height annotations}
For a predicate $p$ on expression nodes, the intended height is
\begin{equation}\label{eq:height}
 h_p(e)=\mathbf1_{p(e)}+\max\bigl(0,\{h_p(c):c\text{ is a child of }e\}\bigr).
\end{equation}
The code's common predicate is recognition of \code{Power[_, _Rational]}.
\code{nestdepths} stores the original expression, all child annotations, and its
computed height in a recursively nested list \src{457--466}.
\code{nestdepthsort} rewrites that tree into groups of equal height, while
\code{maxnestdepth} prunes below the maximum and extracts a frontier
\src{469--498}.

This is a representation measure, not an invariant of the algebraic number.
Algebraically equal expressions can have different heights. The annotation format
also embeds original expression trees at many nodes, which can cause substantial
memory and rewriting overhead on deeply nested inputs. Physical sharing may reduce
actual memory in a particular kernel, so an unconditional measured quadratic-memory
claim would be too strong; the structural duplication risk is nonetheless clear.

The base test uses \code{Depth[expr] === 1}, not an explicit empty-child policy.
A non-atomic expression with no arguments can make the child maximum an empty
maximum, rather than the intended zero in \eqref{eq:height}. This is an edge-domain
risk, not a demonstrated failure on ordinary scalar radical expressions.

The entry test in \code{denest11} looks for arbitrary nested \code{Power} nodes,
whereas the later height computation counts only rational exponents. Moreover,
\code{MemberQ[..., Infinity]} searches levels 1 and deeper, not level 0
\cite{memberq}. A root-level power used as the immediate base of another power
can therefore be missed by that inner test when it has no further power inside.
The gate and the subsequent arithmetic should be brought under one explicit
radical grammar rather than treated as interchangeable tests.

\subsection{The custom factorizer}
\code{fullfactor} begins with \code{FactorList}. It then repeatedly processes
factor/exponent pairs. Numerical rational factors are prime-factorized; sums are
recursively decomposed; a factor that is itself a power has its exponent multiplied
by the outer stored exponent \src{519--548}. \code{combine} merges repeated bases
inside each summand, finds bases common to every summand, subtracts their minimum
exponents, and stores the extracted factors separately \src{554--584}.

\code{Factorc} reconstructs powers, products, and sums using rewrite rules that
recognize the depth and shape of nested lists \src{507--516}. Tagged structures
such as \code{ProductNode}, \code{SumNode}, and \code{PowerNode} would make this
representation considerably easier to validate. As written, ordinary list shape
carries semantic meaning, and the final \code{Flatten[...,][[1]]} idiom assumes the
reconstruction has produced exactly one meaningful expression.

The local \code{j < 5} limit in \code{fullfactor} is an unexplained pass bound.
It may leave some factoring opportunities unexplored. The presence of a remaining
factor does not by itself prove a wrong numerical answer; incomplete factoring
must be distinguished from the genuinely invalid exponent laws below. Also, the
five-pass bound does not bound recursive calls on sums or the larger denesting
search.

\section{Critical correctness defect: the original branch is discarded}
\label{sec:branch}

\subsection{The faulty equality target}
The code sets $A=\alpha^k$, but its candidate filter is
\begin{lstlisting}
Select[possibilities,
  PossibleZeroQ[# - radicand^outerpower] &]
\end{lstlisting}
\src{125--127}. In mathematical notation it tests
\begin{equation}\label{eq:wrongtarget}
 \gamma=A^{1/k}=(\alpha^k)^{1/k},
\end{equation}
not $\gamma=\alpha$.

\begin{proposition}[Exact description of the phase loss]
Let $\alpha\ne0$, write $\alpha=\rho e^{i\theta}$ with
$\theta=\Arg\alpha\in(-\pi,\pi]$, and choose $n\in\Z$ so that
$k\theta-2\pi n\in(-\pi,\pi]$. Then
\begin{equation}\label{eq:phaseloss}
 (\alpha^k)^{1/k}=\alpha e^{-2\pi i n/k}.
\end{equation}
In particular, the reset preserves the target only when $\alpha$ is in the
principal $k$th-root sector $-\pi/k<\Arg\alpha\le\pi/k$.
\end{proposition}
\begin{proof}
The principal logarithm of $\alpha^k$ is
$k\log\rho+i(k\theta-2\pi n)$. Divide it by $k$ and exponentiate.
\end{proof}
Wolfram rational powers use the principal complex root, so this is the relevant
branch convention \cite{power,principalroot}. For $k=2$, the reset can change the
sign. For higher $k$, it can multiply the target by a nontrivial root of unity.

\subsection{Negative real input: the simplest direct-engine example}
Take
\begin{equation}
 \alpha=-\sqrt{3+2\sqrt2}=-(1+\sqrt2).
\end{equation}
The engine chooses $k=2$ and $A=3+2\sqrt2$. With $m=1$, the principal root of $A$
is $1+\sqrt2$, whose minimal polynomial is
\begin{equation}
 P(X)=X^2-2X-1.
\end{equation}
The GCD over $\Q(\sqrt2)$ is $X-1-\sqrt2$. The enumerated candidate set contains
both signs, but the filter selects $1+\sqrt2$, not the negative original input.

A direct reproducer is
\begin{lstlisting}
denest11[-Sqrt[3 + 2 Sqrt[2]], {1}]
\end{lstlisting}
The predicted incorrect result is $1+\sqrt2$. A lone negative sign outside a
radical can remain outside the marker in the public wrapper, so this direct-engine
example alone would not prove that the wrapper always mishandles negatives. The
next two examples address genuinely complex wrapper subproblems.

\subsection{A complex rational power admitted by the wrapper}
Let
\begin{equation}
 z=-1+2i\sqrt2=(1+i\sqrt2)^2.
\end{equation}
The principal square root of $z$ is $1+i\sqrt2$, and
\begin{equation}\label{eq:complexalpha}
 \alpha=z^{3/2}=(1+i\sqrt2)^3=-5+i\sqrt2.
\end{equation}
Its real part is negative, so it is not in the principal square-root sector.
The code's $k=2$ produces
\begin{equation}
 A=\alpha^2=23-10i\sqrt2,\qquad
 A^{1/2}=5-i\sqrt2=-\alpha.
\end{equation}
For $m=1$, the wrong target has minimal polynomial and GCD
\begin{equation}
 P(X)=X^2-10X+27,\qquad G(X)=X-5+i\sqrt2.
\end{equation}
Thus a successful run through the stated trial and filter selects the opposite
value. The relevant public call is
\begin{lstlisting}
DenestRadicals3[(-1 + 2 I Sqrt[2])^(3/2)]
\end{lstlisting}
The mathematics of \eqref{eq:complexalpha}, the square, the polynomial, and the
linear GCD was independently verified. At 30 displayed digits the two values are
\begin{align*}
 \alpha&=-5+1.41421356237309504880168872421\,i,\\
 A^{1/2}&=5-1.41421356237309504880168872421\,i.
\end{align*}
These are not rounding variants; they are exactly negatives of each other.

\subsection{The problem persists even when both outer powers are square roots}
Define
\begin{align}
 u&=\sqrt{-1+2i\sqrt2}=1+i\sqrt2,\\
 v&=\sqrt{-2+2i\sqrt3}=1+i\sqrt3.
\end{align}
Both are principal square roots, yet their product is
\begin{equation}\label{eq:productalpha}
 \alpha=uv=1-\sqrt6+i(\sqrt2+\sqrt3),
\end{equation}
which has negative real part and positive imaginary part. The wrapper's
product-combination rule can send exactly this product to the engine. Squaring gives
\begin{equation}
 A=2-4\sqrt6-i(4\sqrt2+2\sqrt3).
\end{equation}
Its principal square root is $-\alpha$. Independently computed values for the
wrong target are
\begin{align}
 P(X)&=X^4+4X^3+4X^2+48X+144,\\
 G(X)&=X+1-\sqrt6+i(\sqrt2+\sqrt3).
\end{align}
The GCD is linear over a field containing $i\sqrt2$ and $i\sqrt3$. The corresponding
public reproducer is
\begin{lstlisting}
DenestRadicals3[
  Sqrt[-1 + 2 I Sqrt[2]] Sqrt[-2 + 2 I Sqrt[3]]]
\end{lstlisting}
Its intended value is approximately
$-1.4494897427831780982+3.1462643699419723423\,i$; the reset target is its negative.
This exposes the interaction between otherwise value-preserving subproblem grouping
and the engine's invalid assumption about the grouped target.

\subsection{Required correction}
The original target must be passed to \code{trymultiplier} and used in the final
filter. A suitable exact algebraic acceptance test is
\begin{lstlisting}
TrueQ[Quiet[Check[RootReduce[candidate - original] === 0, False]]]
\end{lstlisting}
subject to explicit exact-algebraic input validation and a resource budget.
\code{RootReduce} provides canonical algebraic reduction and can establish exact
equality of algebraic numbers \cite{rootreduce}. Returning the original expression
or an explicit failure when certification is unavailable is safer than accepting a
merely plausible match.

The root-of-unity enumeration should be retained. Replacing the wrong filter by a
correct one does not require guessing the phase numerically: the original target
is already in the candidate orbit under the assumptions proved in
Section~\ref{sec:algebra}.

\section{Other value-changing defects}
\label{sec:othercorrectness}

\subsection{Unconditional flattening of powers in \texttt{fullfactor}}
The branch at source lines 527--531 replaces a factor/exponent pair representing
$(b^r)^s$ by the pair representing $b^{rs}$. The stored outer exponent $s$ can be
rational, because earlier processing has already extracted exponents from radical
nodes. Hence this is not limited to the safe case of raising a power to an integer.

For an unrestricted symbol $z$,
\begin{equation}
 \sqrt{z^2}\ne z
\end{equation}
in general. Setting $z=-1$ gives $1\ne-1$. A symbolic source trace for
\code{Factorc[Sqrt[z^2]]} reaches the exponent multiplication
$2\cdot\tfrac12=1$ and can return $z$. The proposed regression test evaluates the
transformed expression at $z=-1$, so the erroneous transformation cannot be hidden
by comparing two expressions only at a positive sample point.

There is a second branch-law obligation in \code{combine}: extracting a common
factor from a sum and multiplying its exponent by \code{pow} reconstructs a form
like
\begin{equation}
 (cS)^s\longmapsto c^sS^s.
\end{equation}
For $c=S=-1$ and $s=1/2$, the left side is 1 and the right side is $-1$. Appropriate
positivity or argument conditions are needed; the data structure does not record
or check them.

The wrapper's initial \code{Factorc[expr]} means an algebraicity guard on later
radicals does not protect arbitrary symbolic input from these transformations.
Even on a restricted numeric domain, every factoring rewrite used as an equality
must be proved branch-safe or independently checked. These are distinct defects
from the principal-root reset in the engine.

\subsection{Temporary-symbol collisions}
\label{sec:collisions}
\code{Block} localizes values while retaining symbol identities; \code{Module}
creates fresh local symbols \cite{block,module}. That difference is material in a
program that constructs and rewrites symbolic expressions.

\paragraph{The marker \texttt{f} can consume a user function.}
Suppose \code{Global`f} is unassigned and consider
\begin{lstlisting}
DenestRadicals3[f[1], False, Identity]
\end{lstlisting}
The wrapper's \code{Block[{f, ...}]} does not give the marker a new name. Its
\code{Cases[mod, f[_], {0, Infinity}]} therefore recognizes the user's \code{f[1]}
as an internal problem. The callback is applied to 1 and the marker is replaced
by 1. This predicted trace changes a user expression that contains no radical at
all. The same identity collision can interfere with dependency resolution in
all-level mode.

\paragraph{The rationalizer's \texttt{x} can erase a numerator.}
For an unassigned \code{Global`x}, examine
\begin{lstlisting}
RationalizeDenominator[x/(1 + Sqrt[2])]
\end{lstlisting}
The local polynomial indeterminate is also named \code{x}. Source lines 448--450
apply \code{/. x -> 0} to the \emph{entire product} of the numerator and the
polynomial quotient. The user's numerator is therefore replaced by zero as well.
The predicted result is 0 instead of $x(\sqrt2-1)$.

This second example does not require questioning the rationalization identity:
the denominator is a perfectly valid nonzero algebraic number. The error lies in
using the same symbol for user data and a temporary polynomial variable, and in
the scope of the substitution.

Use a private package context and \code{Module} or a fresh \code{Unique} symbol
for markers and polynomial variables. Merely renaming \code{f} to a longer global
name reduces accidental collisions but does not provide lexical hygiene.

\subsection{Loading the file changes a system symbol}
The second source line executes \code{Unprotect[Print]}. No matching restoration
of its previous protection state occurs anywhere in the file. This is a load-time
global side effect, not just a temporary consequence of calling \code{Strad}.
\code{Unprotect} removes the \code{Protected} attribute; ordinary package loading
should not leave unrelated system symbols in a different state \cite{unprotect}.

A better design is a private logging function or an explicit verbosity option.
Expensive debug computations should be placed inside the verbosity branch. Merely
replacing \code{Print} by \code{Null &} suppresses output, but ordinary function
argument evaluation can still compute expressions supplied as log arguments,
including the debug \code{PolynomialGCD} at source line 207 \cite{function}.

\section{Defensive-programming and concrete implementation defects}
\label{sec:defensive}

\subsection{The prime-sorting comparator ignores its second argument}
The comparator at source line 181 is
\begin{lstlisting}
#1[[1]] <= #[[2]] &
\end{lstlisting}
The shorthand \code{#} means \code{#1}, not \code{#2} \cite{function}. On pairs
\code{{prime, exponent}}, this compares the first pair's prime with its own
exponent. It never inspects the second pair.

For example, both comparisons of \code{{2,1}} and \code{{3,1}} return
\code{False}, in either direction. The ordering is not a valid numerical ordering
of the prime factors. Because the next stage assumes an ordered sequence when it
tests ratios of adjacent primes, the typo can alter which multipliers are retained.

The likely intended repair is
\begin{lstlisting}
SortBy[primes, First]
\end{lstlisting}
or \code{#1[[1]] <= #2[[1]] &}. The objection is the wrong slot, not the use of a
non-strict comparison: Wolfram ordering functions can legitimately use such
comparisons \cite{sort}.

\subsection{An empty candidate list is accessed unconditionally}
After filtering candidates, the code handles the case \code{Length[result] > 1},
but not \code{Length[result] == 0}. It then reads \code{result[[1]]}
\src{125--134}. An empty result therefore creates a \code{Part} failure rather
than a controlled unsuccessful trial. The malformed last component may then be
mistaken for a scalar solution by the shape-based success test.

Under ideal exact operations, complete solving, a valid nonzero multiplier, and
a sound equality test, the root-of-unity argument ensures a matching root exists.
That theorem does not excuse the missing guard: invalid inputs, unevaluated solver
results, and heuristic equality decisions can violate those operational
assumptions. A zero forced multiplier is a direct stress case, because it creates
division by $0^{1/k}$ in the candidate construction.

At minimum, require a nonempty validated candidate list before selecting an entry;
otherwise return a typed unsuccessful attempt or \code{Failure}.

\subsection{The equality test is heuristic, not a certificate}
\code{PossibleZeroQ} is documented as using basic symbolic and numerical methods
to suggest whether an expression is zero \cite{possiblezero}. It is appropriate
as an inexpensive screening test, but not as the sole proof of exact branch
identity. In this code it also tests against the wrong target. These are separate
problems: substituting \code{RootReduce} without changing the target would certify
the wrong branch more reliably, not fix the algorithm.

\subsection{Input and intermediate-result validation is missing}
The main functions have broad patterns, but their internals assume explicit exact
algebraic scalars, nonzero multipliers, positive integer root orders at least 2,
well-formed polynomial results, and lists of explicit solutions. These assumptions
are not enforced.

The local \code{trymultiplier} definition only matches \code{Rational[1, _Integer]}
for \code{outerpower} \src{107--108}. If root-order extraction produces 1, a
noninteger, or an unevaluated object, the call can remain unevaluated while later
code still performs positional accesses on \code{attempt}. The source does not
establish that every input passing its initial, broader nested-power test will
produce a valid order.

Lists supplied as forced multipliers are not checked for zero, approximate values,
transcendentals, symbols, or duplicates. Built-ins such as \code{MinimalPolynomial},
\code{FactorInteger}, \code{PolynomialGCD}, and \code{Solve} are followed by indexing
and arithmetic without a consistent \code{Check}, validation, or failure protocol.
An empty forced list itself is not necessarily an error: it can simply yield no
trial and leave the input unchanged. The problem is that these cases have no
explicitly documented status distinction.

\subsection{The ``terms'' extraction is not actually restricted to sums}
The initial-multiplier code uses
\begin{lstlisting}
List @@ Expand[RationalizeDenominator[radicand]]
\end{lstlisting}
\src{224--229}. Replacing a head by \code{List} extracts summands only when the
expanded expression has head \code{Plus}. A \code{Times} expression yields its
factors; a \code{Power} expression yields its base and exponent. An atom is not
converted into the desired singleton list. This makes the heuristic dependent on
an accidental outer head and can propagate unexpected shapes.

A precise decomposition is
\begin{lstlisting}
summands[e_] := If[Head[e] === Plus, List @@ e, {e}]
\end{lstlisting}
with a fresh private definition and separate handling of any desired factor
analysis. The current error mainly affects search quality and robustness; it does
not by itself establish a false equality in the returned expression.

\section{Search limitations and resource risks}
\label{sec:limitations}

\subsection{Equal absolute degrees can hide a useful GCD}
The expensive GCD is attempted only when
\begin{lstlisting}
degree < best || GCD[degree, best] < best
\end{lstlisting}
\src{115}. If a new trial has the same absolute degree as the best known trial,
this test fails. Yet the degree of the GCD over the coefficient field is not
determined by that absolute degree alone.

An explicit construction uses $A=5+2\sqrt6$ from Section~\ref{sec:worked}. The
$m=1$ trial has absolute degree 4 and a degree-2 GCD. Now set
\begin{align}
 \beta&=1+\sqrt2+\sqrt3,\\
 m&=\frac{\beta^2}{A}
   =6-2\sqrt6-2\sqrt2+2\sqrt3>0.
\end{align}
The transformed principal root is $\sqrt{mA}=\beta$. Its minimal polynomial is
\begin{equation}
 P_m(X)=X^4-4X^3-4X^2+16X-8,
\end{equation}
still of absolute degree 4. Over $\Q(\sqrt2,\sqrt3)$, however,
\begin{equation}
 G_m(X)=X-1-\sqrt2-\sqrt3
\end{equation}
is linear. The gate rejects this trial whenever \code{best == 4}, before looking
for the useful GCD. These polynomial facts were independently checked.

This is a counterexample to the \emph{pruning criterion}, not a claim that the
resulting quotient by $\sqrt m$ is the best final denesting, nor a claim that the
automatic generator must propose this particular multiplier. It establishes that
equal absolute degree is not a sound reason to rule out lower-degree GCD information.

\subsection{Partial factorization is intentionally incomplete}
\code{FactorInteger[n, Automatic]} extracts factors that are easy to find; it need
not fully factor $n$ \cite{factorinteger}. The code then deletes composite residual
factors \src{173--181}. Potentially useful prime factors inside a residual composite
are lost. Such a policy can be a reasonable speed tradeoff, but it needs to be
reported as incomplete search, not as evidence that no denesting exists.

Further pruning removes some large, widely separated factors and uses only a
selected prefix in the divisor grid. The discriminant itself is attached to the
chosen polynomial representation; the code supplies no theorem converting this
selection policy into complete coverage of all useful multipliers. In particular,
the GCD-of-degrees branch is not supported by a proved relation between the
corresponding fields.

\subsection{The advertised multiplier cap is not a cap}
Let $r$ distinct primes be retained for \eqref{eq:divisorgrid}. Its divisors have
independent exponents $0,\ldots,k-1$, so their number is exactly
\begin{equation}\label{eq:gridcount}
 k^r.
\end{equation}
The code computes the selected-prime count approximately as
\begin{equation}
 r=\min\!\left(\ell,\max\!\left(5,
 \left\lceil\frac{\log1000}{\log k}\right\rceil\right)\right),
\end{equation}
where $\ell$ is the available prime count after filtering \src{162--194}. A ceiling
usually makes $k^r$ at least the target, and the minimum of five primes can make it
vastly larger. With enough retained primes:
\begin{center}
\begin{tabular}{@{}rrrr@{}}
\toprule
Root order $k$ & Selected primes $r$ & Divisors & After dropping 1\\\midrule
2 & 10 & 1,024 & 1,023\\
3 & 7 & 2,187 & 2,186\\
5 & 5 & 3,125 & 3,124\\
10 & 5 & 100,000 & 99,999\\
100 & 5 & 10,000,000,000 & 9,999,999,999\\
\bottomrule
\end{tabular}
\end{center}
These are exact combinatorial counts, not measured trial counts. The union with
additional discovered primes can add still more candidates. The variable
\code{maxprimes = 14} is declared but never used.

A genuine bound for the grid alone would require $k^r-1\le B$, with the remaining
budget reserved for extra primes and existing proposals. Better still, generate
exponent vectors lazily and stop after $B$ accepted proposals. Calling
\code{Divisors} and then truncating its output is not a memory-safe implementation
of a cap, because the large list has already been materialized.

\subsection{The subset pass adds another exponential mechanism}
If $q$ promising initial multipliers are selected, the \code{Distribute} construction
can form up to $2^q$ subset products before duplicates are removed \src{402--420}.
The divisor-grid cap, even if repaired, would not bound this separate expansion.
The rational powers and coefficient heights can also increase the cost of a single
minimal-polynomial, discriminant, or GCD computation enormously.

There is no global visited set of canonical multiplier keys, so the same algebraic
multiplier can be proposed through different frames or in different syntactic
forms. There is no explicit budget for total trials, total queued proposals,
polynomial degree, coefficient size, time, or memory.

\subsection{Termination and failure status need a real contract}
The outer \code{For} loop compares its index with the current length of a worklist
that can grow while it is being traversed. The source provides no monotone measure
proving that this growth stops. The history often improves and the heuristics may
work well in practice, but that is not a termination proof.

It would be inaccurate to say every \code{//.} loop is completely unbounded:
\code{ReplaceRepeated} has a documented default \code{MaxIterations} limit of
65,536 \cite{replacerepeated}. That built-in bound is not a global search budget,
and hitting it is not a meaningful ``no denesting exists'' outcome. Kernel
recursion or evaluation limits are similarly not substitutes for a deliberate
algorithmic status.

The final messages distinguish mainly a changed expression from ``no progress''
\src{434--440}. They do not distinguish search exhaustion, incomplete
factorization, skipped equal-degree trials, failed algebraic operations,
certification failure, and resource exhaustion. These are different outcomes and
should be returned as data.

\subsection{Sorting and simplification have hidden costs}
\code{ComplexitySort} prioritizes \code{LeafCount} and uses \code{Abs} to break ties
\src{500--504}. This is a heuristic structural order, not a measure of algebraic
field complexity or denesting quality. If exact comparisons remain unresolved,
the ordering function may not provide the intended numerical ordering; forcing
such comparisons can itself be costly. The initial sort by real part has a
similar limitation outside the intended explicit-number domain.

Repeated \code{Expand}, \code{Together}, \code{Simplify}, discriminant factoring,
and tree reconstruction add overhead before a useful multiplier is found. No
benchmarks are supplied, so this report does not rank their empirical contribution.
Caching algebraic computations and counting work by stage would make such an
assessment possible.

\section{Consolidated findings and repair priorities}
\label{sec:findings}

Severity here describes consequences, not measured frequency. ``High'' includes
silent value changes and uncontrolled resource growth; heuristic incompleteness is
kept separate from soundness.
\begin{longtable}{@{}>{\raggedright\arraybackslash}p{0.07\textwidth}>{\raggedright\arraybackslash}p{0.12\textwidth}>{\raggedright\arraybackslash}p{0.48\textwidth}>{\raggedright\arraybackslash}p{0.24\textwidth}@{}}
\toprule
ID & Priority & Finding & Evidence and location\\\midrule
\endhead
B1 & Critical & Candidate selection targets $(\alpha^k)^{1/k}$ rather than
$\alpha$. & Algebraic counterexamples; 90--94, 125--127.\\
B2 & Critical & Custom factorization uses unrestricted complex power laws. &
Source rule and branch counterexamples; 527--543, 575--583.\\
B3 & High & Non-fresh marker \code{f} can consume caller data. & Predicted source
trace; 6, 55--73.\\
B4 & High & Rationalizer's \code{x -> 0} can erase a numerator named \code{x}. &
Predicted source trace; 445--450.\\
B5 & High & Empty candidate list is indexed as a success. & Unconditional access;
125--134.\\
B6 & High & Multiplier ``cap'' admits exponential candidate lists. & Exact count
$k^r$; 162--199.\\
B7 & High & No global work, time, memory, or queued-candidate limit. & Growing
worklist and subset generation; 261--430.\\
B8 & Medium & Prime comparator compares a pair with itself. & Wrong pure-function
slot; 181.\\
B9 & Medium & Loading removes \code{Print} protection without restoration. &
Direct load-time side effect; 2.\\
B10 & High & No consistent validation or failed-operation protocol. & Broad entry
patterns, narrow trial pattern, positional assumptions; 79--134, 214--215.\\
H1 & Medium & Equal-degree gate can skip a linear GCD. & Explicit algebraic
counterexample; 115.\\
H2 & Medium & Partial-factorization residuals are discarded. & Documented built-in
behavior and source policy; 173--181.\\
H3 & Medium & Discriminant and degree-GCD searches have no coverage proof. &
Heuristic construction; 169--207.\\
H4 & Medium & Initial ``terms'' may actually be factors or power arguments. &
Head replacement; 224--229.\\
H5 & Medium & No certified radical-form or complexity improvement criterion. &
Success selection and replacement; 120--134, 309--324.\\
M1 & Lower & Opaque list schemas, fixed factoring passes, and unused variables. &
Maintenance and auditability; 95--103, 163, 253--267, 522.\\
\bottomrule
\end{longtable}

Several source variables are unused or only diagnostic: \code{extrapower} is
printed but not consumed, \code{maxprimes} is never applied, \code{expected} is
unused inside a trial, \code{done} is unused inside the multiplier generator, and
\code{best} in the extra-multiplier block is calculated but not used. The many
commented-out alternatives and version-numbered function names suggest an
experimental script, but they do not identify an author, development date, or
historical correctness guarantee.

\section{A safer design that preserves the useful mathematics}
\label{sec:repair}

\subsection{Separate proposal generation from proof}
A robust implementation can retain the multiplier heuristic while treating it as
an untrusted proposer. The acceptance layer should be much smaller and easier to
audit. It must retain the original target, verify every proposed result exactly,
and enforce the chosen output language and improvement criterion.

The required invariants are:
\begin{enumerate}
\item The semantic target never changes. A powered radicand is auxiliary data;
it does not replace the original input for equality testing.
\item Every transformation used as an equality either has proved branch
conditions or passes exact verification against its own pre-transformation value.
\item A candidate is called ``denested'' only after exact equality and the output
representation criterion both pass.
\end{enumerate}
These checks belong around the entire wrapper, not only around one internal trial.
For exact numeric inputs, a final check against the original \code{expr} can catch
preprocessing errors. For symbolic inputs, either require a clearly specified
assumption system with proved transformations or reject them explicitly.

\subsection{Use exact algebraic verification conservatively}
For explicit exact algebraic numbers, one can screen cheaply and then certify by
reducing the difference to zero. A failed or timed-out verification should not be
reinterpreted as equality. The strengthened acceptance condition is
\begin{equation}
 \operatorname{RootReduce}(\gamma-\alpha)=0,
\end{equation}
with the actual original target $\alpha$.

An acceptance layer should also reject \code{Indeterminate}, infinities, unresolved
variables, approximate numbers, \code{ConditionalExpression} objects outside the
contract, and failed solver calls. \code{TrueQ} is useful where a total Boolean
decision is required, but replacing every test by \code{TrueQ} does not turn an
undecidable or unevaluated expression into a proof; it merely makes failure
conservative.

\subsection{Choose a denesting objective explicitly}
One possible objective is a lexicographic score built from radical nesting height,
number of rational-power nodes, and expression size. It should be accompanied by
an output-language predicate that disallows \code{Root} or \code{AlgebraicNumber}
when the task specifically requires radicals. Otherwise a canonical algebraic
number representation can win a syntactic score without satisfying the intended
mathematical goal.

A candidate may be exactly equal but not improved. One suitable status is
\code{"EqualButNotImproved"}. It should not be described as \code{"Denested"}.
When several certified expressions are available, compare their representation
scores instead of taking the first structurally distinct expression returned by
\code{Solve}.

\subsection{Make state and resources explicit}
Use associations or another tagged record type for attempts and work items. For
example, distinguish
\begin{lstlisting}
<|"Status" -> "NoLowerDegreeGCD", "Degree" -> d|>
<|"Status" -> "CertifiedCandidates", "Candidates" -> values|>
Failure["NoCertifiedCandidate", <|"Multiplier" -> m|>]
Failure["TimeLimit", <|"TrialsCompleted" -> count|>]
\end{lstlisting}
A deterministic queue, a visited set, lazy candidate generation, and an actual
trial budget should replace the implicit growth policy. Include the root order,
coefficient-field context, and any other relevant trial state in cache keys.
Canonicalization can identify algebraically identical multipliers, but it also has
a cost and must be budgeted.

For pure rational multipliers, reducing prime valuations modulo $k$ can sometimes
organize equivalent search classes. Such an optimization needs a proof appropriate
to the trial and branch conventions; it should not be silently inferred from the
current exponent-reduction rules.

\subsection{A guarded one-multiplier step}
The archive includes \code{safer_core.wl}. It is a deliberately limited reference
component, not a patched replacement for all 585 source lines. It validates exact
algebraic input, validates $k\ge2$ and $m\ne0$, computes the same GCD idea without
absolute-degree pruning, enumerates roots of unity, certifies against the original
target, handles an empty candidate set, and returns structured results under a time
limit. Its key acceptance block is
\begin{lstlisting}
roots = t /. rules;
candidates = Flatten[Outer[Times,
    roots/multiplier^(1/k),
    Exp[2 Pi I Range[0, k - 1]/k]]];
accepted = Select[candidates,
  exactAlgebraicQ[#] && ExactAlgebraicEqualQ[#, target] &];
If[accepted === {},
  Return[Failure["NoCertifiedCandidate", <|
    "Multiplier" -> multiplier, "GCDDegree" -> dg|>]]];
\end{lstlisting}
The complete component is supplied as source, with an explicit warning that it was
not executed in a Wolfram kernel during this audit. It certifies equality, not
radical denesting; it can still return a representation that a higher-level
radical-only policy would reject. It also does not implement a complete search or
a global memory budget. Keeping these boundaries explicit is preferable to calling
a small illustrative repair a production-ready denester.

\section{Verification performed and proposed regression testing}
\label{sec:tests}

\subsection{Independently executed checks}
The included \code{verify_algebra.py} script was executed successfully. It uses
SymPy for exact minimal polynomials, polynomial GCDs over explicit algebraic
extensions, polynomial division, and identity checks; it uses mpmath at 80 decimal
digits to check the principal-branch examples. Its text output and JSON output
are included.

The checks cover the positive worked example, the negative branch example, the
complex $3/2$ power, the merged square-root product, the equal-absolute-degree
pruning counterexample, three denominator-rationalization examples, the two
invalid power identities, the malformed comparator's logic, and exact
multiplier-grid counts. Numerical results supplement the exact arguments rather
than replacing them.

These checks validate the algebra used in the audit. They do not establish the
complete behavior of \code{FactorList}, \code{Solve}, evaluation scoping, expression
canonicalization, or the mutable search stack in a specific Wolfram release.

\subsection{Wolfram tests supplied but not executed}
The file \code{regression_tests.wlt} contains expected-behavior tests for the
original source and a few tests of the guarded component. Several expectations
are deliberately preservation requirements that the legacy code is predicted to
violate; they are not presented as a passing legacy suite.

The suite checks that loading preserves \code{Print} attributes, that the prime
comparator compares the two prime values, that negative and complex branches are
preserved, that a symbolic square is not flattened incorrectly, that a caller's
\code{f} is not consumed, and that a numerator named \code{x} survives
rationalization. Positive examples and all-level value preservation are included
as controls. The guarded component is tested for rejection of a zero multiplier
and invalid root order and for preservation of negative and complex targets.

Because loading the original file changes \code{Print}, these tests should be run
in a disposable, fresh kernel. The suite records and restores the prior attribute
list after loading. A typical invocation from the archive directory is
\begin{lstlisting}[language={}]
wolframscript -code 'TestReport["regression_tests.wlt"]'
\end{lstlisting}
Record \code{$Version}, \code{$SystemID}, \code{$Assumptions}, relevant evaluation
limits, and the test report alongside results. A failure caused by automatic
pre-evaluation must be interpreted against the actual input reaching the function,
not assumed to invalidate the mathematical branch counterexample.

\subsection{Tests needed beyond the supplied regression file}
A broader suite should cover conjugate inputs, branch-cut boundary values,
negative rational exponents, mixed root orders, integer powers containing radicals,
zero and singular expressions, symbolic and approximate inputs, invalid callbacks,
and repeated or empty multiplier lists. It should test every public entry point,
not only \code{denest11}.

Useful properties include exact value preservation, no leakage of private markers,
no mutation of unrelated definitions or attributes, stability under renaming
unrelated symbols, and explicit resource-limit results. Idempotence is desirable
for a normalizer, but should not be imposed as a theorem on an intentionally
budget-limited heuristic without specifying how the budget is used.

A performance harness should separately record preprocessing time, minimal-
polynomial time, GCD time, factorization time, number of unique trials, peak queued
candidates, output score, and exact-verification time. No performance improvement
percentage is claimed here because none was measured.

\section{Conclusion}
The program contains a substantive algebraic search idea. Its use of a minimal
polynomial together with an algebraic-coefficient GCD can reveal a simpler relation,
and its enumeration of roots of unity is a mathematically appropriate way to
recover possible branches. Its denominator-rationalization formula is also valid
for nonzero algebraic denominators.

Those strengths do not make the implementation a reliable general denester in its
current form. The original target is not preserved during candidate selection;
the custom factorizer performs branch-sensitive rewrites without conditions;
non-fresh temporary symbols can alter caller expressions; and several error paths
and resource limits are missing. The sorting typo and the ineffective multiplier
cap are concrete implementation defects, while partial factorization and
unproved degree-based pruning explain why a failed search says little about the
existence of a denesting.

The first repair priority is correctness, not a larger multiplier search: retain
the original expression, use branch-safe preprocessing, make temporary symbols
fresh, validate operations, and require exact equality before accepting anything.
Only then should the search heuristics be optimized and benchmarked. With that
separation, the current multiplier strategy can be retained as one useful proposal
engine without being asked to serve as its own proof of correctness or completeness.

\begin{thebibliography}{99}
\addcontentsline{toc}{section}{References}
\bibitem{source} Supplied attachment, \emph{Pasted text.txt}. Original bytes preserved
as \texttt{original.wl}; SHA-256 stated on the title page. Source reproduced in
Appendix~\ref{app:source}. No author or version metadata supplied.
\bibitem{minpoly} Wolfram Research, \emph{MinimalPolynomial}, Wolfram Language
Documentation. \url{https://reference.wolfram.com/language/ref/MinimalPolynomial.html}.
\bibitem{polygcd} Wolfram Research, \emph{PolynomialGCD}, Wolfram Language
Documentation. \url{https://reference.wolfram.com/language/ref/PolynomialGCD.html}.
\bibitem{rootreduce} Wolfram Research, \emph{RootReduce}, Wolfram Language
Documentation. \url{https://reference.wolfram.com/language/ref/RootReduce.html}.
\bibitem{power} Wolfram Research, \emph{Power}, Wolfram Language Documentation.
\url{https://reference.wolfram.com/language/ref/Power.html}.
\bibitem{principalroot} Wolfram MathWorld, \emph{Principal Root}.
\url{https://mathworld.wolfram.com/PrincipalRoot.html}.
\bibitem{block} Wolfram Research, \emph{Block}, Wolfram Language Documentation.
\url{https://reference.wolfram.com/language/ref/Block.html}.
\bibitem{module} Wolfram Research, \emph{Module}, Wolfram Language Documentation.
\url{https://reference.wolfram.com/language/ref/Module.html}.
\bibitem{possiblezero} Wolfram Research, \emph{PossibleZeroQ}, Wolfram Language
Documentation. \url{https://reference.wolfram.com/language/ref/PossibleZeroQ.html}.
\bibitem{factorinteger} Wolfram Research, \emph{FactorInteger}, Wolfram Language
Documentation. \url{https://reference.wolfram.com/language/ref/FactorInteger.html}.
\bibitem{sort} Wolfram Research, \emph{Sort}, Wolfram Language Documentation.
\url{https://reference.wolfram.com/language/ref/Sort.html}.
\bibitem{function} Wolfram Research, \emph{Function}, Wolfram Language Documentation.
\url{https://reference.wolfram.com/language/ref/Function.html}.
\bibitem{replacerepeated} Wolfram Research, \emph{ReplaceRepeated}, Wolfram Language
Documentation. \url{https://reference.wolfram.com/language/ref/ReplaceRepeated.html}.
\bibitem{memberq} Wolfram Research, \emph{MemberQ}, Wolfram Language Documentation.
\url{https://reference.wolfram.com/language/ref/MemberQ.html}.
\bibitem{complexexpand} Wolfram Research, \emph{ComplexExpand}, Wolfram Language
Documentation. \url{https://reference.wolfram.com/language/ref/ComplexExpand.html}.
\bibitem{unprotect} Wolfram Research, \emph{Unprotect}, Wolfram Language Documentation.
\url{https://reference.wolfram.com/language/ref/Unprotect.html}.
\bibitem{factorlist} Wolfram Research, \emph{FactorList}, Wolfram Language
Documentation. \url{https://reference.wolfram.com/language/ref/FactorList.html}.
\end{thebibliography}
All external documentation above was consulted on 5 September 2026. The source
listing and original analysis, rather than an external denesting algorithm or an
unseen specification, are the basis of the audit.

\clearpage
\appendix
\section{Complete supplied source, with physical line numbers}
\label{app:source}
The following listing is the supplied program, not a corrected version. Line
numbers refer to the original physical lines, including blank lines. Long lines
may wrap in the PDF without receiving a new line number. The file in the archive
preserves the original CRLF byte stream and its SHA-256 fingerprint.
\lstinputlisting[numbers=left,firstnumber=1,basicstyle=\ttfamily\scriptsize,
 xleftmargin=22pt,numbersep=7pt]{original.wl}

\clearpage
\section{Archive contents and reproducibility}
The requested deliverables are \code{article.tex} and \code{article.pdf}. The
additional files make the reasoning and proposed runtime checks auditable:
\begin{longtable}{@{}p{0.35\textwidth}p{0.61\textwidth}@{}}
\toprule
File & Purpose\\\midrule
\endhead
\code{original.wl} & Original attachment, byte-for-byte.\\
\code{verify_algebra.py} & Executed independent algebra and branch checks.\\
\code{verification_results.txt} & Captured output of that execution.\\
\code{verification_results.json} & Machine-readable independent check results.\\
\code{safer_core.wl} & Statically reviewed, unexecuted guarded trial component.\\
\code{regression_tests.wlt} & Unexecuted Wolfram expected-behavior tests.\\
\code{README.md} & Build and test instructions, plus the execution boundary.\\
\code{SHA256SUMS.txt} & Checksums of archive content other than the checksum file.\\
\bottomrule
\end{longtable}
Rebuild the article from this directory with
\begin{lstlisting}[language={}]
pdflatex -interaction=nonstopmode -halt-on-error article.tex
pdflatex -interaction=nonstopmode -halt-on-error article.tex
\end{lstlisting}
The original source file is required alongside the LaTeX file because the full
listing is included with \code{\lstinputlisting}. No external image or font file is
needed. Re-run the independent checks with
\begin{lstlisting}[language={}]
python verify_algebra.py
\end{lstlisting}
The Python script requires SymPy and mpmath. It does not require, contact, or
substitute for a Wolfram kernel. Re-running the Wolfram tests is a separate step.
\end{document}
